\documentclass[11pt,a4paper]{article}

\usepackage[utf8]{inputenc}
\usepackage[T1]{fontenc}
\usepackage[margin=2.5cm]{geometry}
\usepackage{amsmath,amssymb}
\usepackage{booktabs}
\usepackage[hidelinks]{hyperref}

\title{Graph Neural Networks for Molecular Property Prediction:\\
       Expressiveness, Failure Modes, and What Actually Transfers}
\author{Prepared with LaTeXEditor}
\date{July 2026}

\begin{document}
\maketitle

\begin{abstract}
Molecules are graphs, and the natural way to learn functions of them is with
architectures that respect that structure. This paper reviews message-passing
neural networks for molecular property prediction, establishes what they can
and cannot represent, and examines three failure modes --- over-smoothing,
over-squashing, and the expressiveness ceiling imposed by the
Weisfeiler--Leman test --- that are frequently discussed separately but which
interact. We then turn to a less theoretical question: whether reported gains
transfer. A recurring finding is that carefully tuned baselines close much of
the gap to sophisticated architectures on standard benchmarks, and that
scaffold-split evaluation, which tests generalisation to structurally novel
compounds, reverses several published rankings. We argue that the field's
central difficulty is not architectural but distributional, and that progress
is limited more by the size and coverage of labelled chemical data than by the
expressiveness of the models applied to it.
\end{abstract}

\tableofcontents

\section{Introduction}
\label{sec:intro}

The properties a chemist cares about --- solubility, toxicity, binding
affinity, synthetic accessibility --- are functions of molecular structure.
Predicting them computationally has a long history, and until recently the
dominant approach was to hand-design a fixed-length vector of structural
features, a \emph{fingerprint}, and fit a conventional model to it.

Fingerprints work well and remain competitive. Their limitation is that the
features are fixed in advance: the representation cannot adapt to the property
being predicted. Learning the representation instead was proposed by
\cite{duvenaud2015}, who replaced the fixed hashing of circular fingerprints
with a differentiable analogue, and this is the direct ancestor of the
message-passing networks now standard.

The appeal is obvious. A molecule is naturally a graph --- atoms are nodes,
bonds are edges --- and a model operating directly on that graph respects the
symmetries of the object. Permuting the atom indices should not change a
predicted boiling point, and a graph network guarantees this by construction
rather than learning it from data.

Whether the appeal has been realised is a more complicated question, and it is
the subject of this paper. Sections~\ref{sec:mpnn} and \ref{sec:expressive}
establish the framework and its representational limits.
Section~\ref{sec:failure} examines three failure modes.
Sections~\ref{sec:geometry} and \ref{sec:transformers} cover the two main
lines of response. Section~\ref{sec:evaluation} then argues that the
evaluation practices of the field have systematically overstated progress.

\section{Message passing}
\label{sec:mpnn}

Most graph networks are instances of a single framework \cite{gilmer2017}.
Each node $v$ carries a state $h_v$, updated over $T$ rounds by aggregating
messages from its neighbours $\mathcal{N}(v)$:
\begin{align}
  m_v^{(t+1)} &= \sum_{u \in \mathcal{N}(v)}
                 M_t\big(h_v^{(t)}, h_u^{(t)}, e_{uv}\big),
  \label{eq:message} \\
  h_v^{(t+1)} &= U_t\big(h_v^{(t)}, m_v^{(t+1)}\big),
  \label{eq:update}
\end{align}
where $e_{uv}$ are edge features, $M_t$ is a message function and $U_t$ an
update function. A graph-level prediction is obtained by pooling:
\begin{equation}
  \hat{y} = R\big(\{ h_v^{(T)} : v \in G \}\big),
\end{equation}
with $R$ a permutation-invariant readout, typically a sum or mean followed by
a feed-forward network.

The design space is the choice of $M$, $U$ and $R$. Graph convolutional
networks \cite{kipf2017} use a degree-normalised mean; GraphSAGE
\cite{hamilton2017} samples a fixed-size neighbourhood and concatenates rather
than sums; graph attention networks \cite{velickovic2018} weight neighbours by
a learned attention coefficient,
\begin{equation}
  \alpha_{vu} = \frac{\exp\big(\sigma(a^{\top}[Wh_v \| Wh_u])\big)}
                      {\sum_{k \in \mathcal{N}(v)}
                       \exp\big(\sigma(a^{\top}[Wh_v \| Wh_k])\big)}.
\end{equation}

\paragraph{The receptive field.}
After $T$ rounds a node's state depends on its $T$-hop neighbourhood, and no
further. This is the graph analogue of a convolutional receptive field, and it
has an immediate chemical consequence: a property determined by the
relationship between two distant functional groups cannot be represented
unless $T$ exceeds the distance between them. Typical implementations use
$T = 3$ to $5$, which covers local chemistry well and long-range interactions
not at all.

\section{Expressiveness}
\label{sec:expressive}

There is a precise characterisation of what message passing can distinguish
\cite{xu2019, morris2019}: no message-passing network can distinguish two
graphs that the one-dimensional Weisfeiler--Leman (1-WL) colour refinement
test fails to distinguish.

The argument is direct. 1-WL iteratively refines a colouring by hashing each
node's colour together with the multiset of its neighbours' colours ---
structurally the same operation as Equations~\ref{eq:message}
and~\ref{eq:update}. A network can therefore be no more discriminative than
the test, and is exactly as discriminative when its aggregation is injective.

Graph isomorphism networks \cite{xu2019} achieve this bound by using a sum
aggregator with a multi-layer perceptron,
\begin{equation}
  h_v^{(t+1)} = \mathrm{MLP}\Big( (1 + \epsilon) h_v^{(t)} +
                \sum_{u \in \mathcal{N}(v)} h_u^{(t)} \Big),
\end{equation}
since sum is injective over multisets in a way that mean and max are not: mean
cannot distinguish a node with two identical neighbours from one with one, and
max discards multiplicity entirely.

\paragraph{Does the ceiling matter chemically?}
The standard illustration is that 1-WL cannot distinguish certain regular
graphs, and the standard chemical instance is that it cannot count cycles.
Since aromaticity is a ring property with strong effects on reactivity, this
looks serious. In practice it is less so than it appears, because molecular
graphs are annotated: atom features typically include ring membership and
aromaticity flags computed by a cheminformatics toolkit. The information that
1-WL cannot derive is supplied as input.

This is worth stating plainly, because a good deal of architectural work is
motivated by the WL bound while the empirical gain from exceeding it is
frequently small. The bound constrains what the network can \emph{compute};
it says nothing about what it must compute when the features already encode
the answer.

\begin{table}[htbp]
  \centering
  \caption{Aggregation functions and their injectivity over multisets.}
  \label{tab:aggregation}
  \begin{tabular}{llll}
    \toprule
    \textbf{Aggregator} & \textbf{Injective?} & \textbf{Distinguishes} &
      \textbf{Typical use} \\
    \midrule
    Sum \cite{xu2019}       & yes & multiset          & expressiveness-critical \\
    Mean \cite{kipf2017}    & no  & distribution      & node classification \\
    Max \cite{hamilton2017} & no  & underlying set    & skeleton detection \\
    Attention \cite{velickovic2018} & no & weighted distribution & heterogeneous graphs \\
    \bottomrule
  \end{tabular}
\end{table}

\section{Three failure modes}
\label{sec:failure}

\subsection{Over-smoothing}

Stacking message-passing layers makes node representations converge. Each
round averages a node with its neighbours, and repeated averaging is a
diffusion process whose fixed point is constant over any connected component
\cite{li2018}. After enough rounds every node carries the same vector and all
structural information is lost.

The practical effect is that graph networks are shallow. Where image
classifiers went from eight layers to over a hundred and improved throughout,
graph networks typically peak at three to five and degrade after. Residual
connections, jumping knowledge and explicit normalisation all mitigate the
problem without removing it.

\subsection{Over-squashing}

The second failure is nearly the opposite \cite{alon2021}. Information from an
exponentially growing $T$-hop neighbourhood must be compressed into a
fixed-width node vector. Where the graph has a bottleneck --- a bridge edge
through which many paths must pass --- the messages crossing it are squashed,
and long-range signal is lost regardless of depth.

The two failures place the architecture in a vice. Adding layers to reach
distant nodes worsens over-smoothing; keeping the model shallow leaves distant
nodes unreachable. Molecules are relatively small and often tree-like, which
limits the damage, but for the long-thin molecules common in medicinal
chemistry the bottleneck is real.

\subsection{Depth without benefit}

Together these produce the field's most conspicuous empirical regularity: on
most molecular benchmarks, a well-tuned three-layer network is as good as
anything deeper. The result is stable across datasets and architectures, and
it should be read as a statement about the mechanism rather than about tuning
effort.

\section{Adding geometry}
\label{sec:geometry}

A molecular graph omits something a chemist would consider essential:
three-dimensional structure. Two stereoisomers have identical graphs and
different properties. Any purely topological model is blind to the difference.

Geometric networks take atomic coordinates as input. The difficulty is that
predictions must be invariant to rotation and translation --- rotating a
molecule does not change its energy --- and naive use of coordinates violates
this. SchNet \cite{schutt2017} obtains invariance by using only interatomic
distances, expanded in a radial basis:
\begin{equation}
  h_v^{(t+1)} = h_v^{(t)} + \sum_{u \ne v}
    W\big(\| r_v - r_u \|\big) \odot h_u^{(t)}.
\end{equation}
Distances are invariant, so the model is. But distances alone do not determine
geometry: two conformations can share all pairwise distances between bonded
atoms and differ in torsion. DimeNet \cite{gasteiger2020} adds angular
information via directional messages, and E($n$)-equivariant networks
\cite{satorras2021} go further, producing outputs that transform correctly
under rotation rather than being merely invariant.

\paragraph{The conformer problem.}
Geometric models require coordinates, and coordinates require a conformer.
Molecules are flexible; the relevant conformer is often the bound one, which
is unknown. Generating conformers computationally introduces error that is
frequently larger than the difference between architectures being compared.
This is a data problem masquerading as a modelling problem, and it accounts
for much of the inconsistency in reported geometric-model performance.

\section{Graph transformers}
\label{sec:transformers}

The response to over-squashing that has gained most ground is to abandon
locality. If every node attends to every other, bottlenecks are irrelevant and
the receptive field is the whole graph after one layer \cite{ying2021}.

The immediate objection is that this discards the structure entirely --- a
transformer over nodes has no notion of which are bonded. Graph transformers
reinject it as a bias on the attention scores, typically a learned function of
shortest-path distance:
\begin{equation}
  A_{ij} = \frac{(h_i W^Q)(h_j W^K)^{\top}}{\sqrt{d}} + b_{\phi(i,j)},
\end{equation}
where $\phi(i,j)$ is the distance between nodes. Hybrid designs
\cite{rampasek2022} run message passing and global attention in parallel and
combine them, on the reasoning that local chemistry is best handled by the
former and long-range effects by the latter.

Molecules are small enough that the quadratic cost of attention is not
prohibitive --- a hundred atoms gives a $100 \times 100$ matrix --- so the
trade that dominates the sequence literature does not bind here. The
constraint is data rather than compute.

\section{Evaluation, and why the rankings move}
\label{sec:evaluation}

MoleculeNet \cite{wu2018} standardised evaluation and is the most-used
benchmark suite. Its adoption has been valuable and has also entrenched
several practices that flatter new methods.

\paragraph{Split choice dominates.}
A random split places structurally similar molecules in both train and test,
so the task rewards interpolation. A scaffold split groups by core structure,
so the test set is structurally novel --- the situation an actual discovery
programme faces. Scaffold performance is invariably worse, often dramatically,
and the ranking of methods under the two splits is not the same. Results
reported under random splits should be treated as upper bounds that do not
describe the intended use case.

\paragraph{Baselines are undertuned.}
Repeated reproduction studies have found that gradient-boosted trees on
conventional fingerprints, given tuning effort comparable to that lavished on
the proposed model, match or beat graph networks on a majority of MoleculeNet
tasks. The datasets are small --- many have a few thousand labelled compounds
--- and in that regime a learned representation has little advantage over a
good fixed one.

\paragraph{Datasets are small and biased.}
Labelled molecular data is expensive: each point is an experiment. The
datasets are correspondingly small and drawn from historical medicinal
chemistry programmes, so they cover the regions of chemical space that
chemists have already explored. A model that generalises within that region
may be useless outside it, and the benchmark cannot detect the difference.

\begin{table}[htbp]
  \centering
  \caption{Effect of split choice on reported performance. Values are
           indicative of the pattern reported across the literature rather
           than a single measurement.}
  \label{tab:splits}
  \begin{tabular}{lll}
    \toprule
    \textbf{Split} & \textbf{Tests} & \textbf{Reported performance} \\
    \midrule
    Random             & interpolation within known chemistry & highest \\
    Scaffold           & generalisation to novel cores        & substantially lower \\
    Time-based         & prospective prediction               & lowest \\
    \bottomrule
  \end{tabular}
\end{table}

\section{Where the approach has worked}
\label{sec:success}

The critical tone above should not obscure genuine successes, which share a
common shape: large screening libraries, a clear binary label, and use of the
model as a filter rather than an oracle.

The best-known is the identification of halicin as an antibacterial
\cite{stokes2020}. A graph network trained on a growth-inhibition screen was
used to rank a large library, and a highly ranked compound with no structural
similarity to known antibiotics was confirmed experimentally. The model did
not predict a property precisely; it reduced a hundred million candidates to a
few hundred worth testing, which is exactly what a noisy but cheap predictor
is good for.

Structure prediction \cite{jumper2021} is the other landmark, and although the
architecture is not a message-passing network in the sense of
Section~\ref{sec:mpnn}, it shares the essential feature of operating on a
structured representation with the right symmetries. Its success rests on
something the molecular property field lacks: a large, curated, consistent
body of experimental structures accumulated over decades.

\section{Pretraining and transfer}
\label{sec:pretraining}

The size of labelled molecular datasets is the binding constraint identified
in Section~\ref{sec:evaluation}, and the obvious response is the one that
worked in language: pretrain on abundant unlabelled data, fine-tune on the
scarce labelled task. Unlabelled molecules are genuinely abundant --- public
databases hold hundreds of millions of structures --- so the ingredients are
present.

Several pretext tasks have been proposed. Node-level objectives mask an atom
and predict its identity from context, directly analogous to masked language
modelling. Edge-level objectives predict whether a bond exists between two
atoms. Graph-level objectives predict cheap computed descriptors, or bring
representations of two conformers of the same molecule together while pushing
apart those of different molecules.

The results have been mixed in an instructive way. Node-level masking in
particular can \emph{hurt} downstream performance, a phenomenon termed
negative transfer. The likely explanation is that predicting a masked atom is
nearly trivial given valence rules --- a carbon with four single bonds to
carbons is overwhelmingly a carbon --- so the objective is satisfied by
learning local chemistry the model would learn anyway, while the
representation is shaped by a task with no bearing on the property of
interest.

This contrasts sharply with language, where masked-token prediction is
genuinely hard and requires broad understanding. The lesson is that a pretext
task must be difficult to be useful, and difficulty is a property of the
domain rather than of the objective's form. Objectives that predict computed
quantum properties, which are expensive to calculate but cheap to obtain in
bulk from simulation, transfer more reliably --- they are hard, and their
difficulty is of the right kind.

\section{Uncertainty and calibration}
\label{sec:uncertainty}

A property predictor used to prioritise experiments is only as useful as its
uncertainty estimates. Ranking a library by predicted activity is worthwhile
only if the model can also indicate where it is guessing, since the compounds
most worth testing are frequently those the model has least reason to be
confident about.

Neural networks are notoriously overconfident, and graph networks are no
exception. The standard remedies transfer directly: deep ensembles, which
train several models with different initialisations and use the spread of
their predictions; Monte Carlo dropout, which is cheaper and generally worse;
and post-hoc calibration on a held-out set.

The molecular setting adds a specific difficulty. The relevant uncertainty is
usually \emph{epistemic} --- the model has not seen anything like this
compound --- rather than aleatoric noise in the assay. Epistemic uncertainty
should grow with distance from the training distribution, and the natural
measure of that distance is chemical similarity. Methods that combine a
learned predictor with an explicit similarity-to-training-set term are
consistently better calibrated on scaffold splits than methods relying on the
network's own confidence, precisely because the network has no representation
of what it has not seen.

\paragraph{Applicability domain.}
Classical cheminformatics has a name for this: the applicability domain of a
model, the region of chemical space in which its predictions can be trusted.
The concept predates deep learning and is arguably more important now, since
larger models extrapolate more confidently and no less wrongly. Reporting a
prediction without an applicability judgement is an omission the older
literature would not have permitted.

\section{Interpretation}
\label{sec:interpretation}

A medicinal chemist asked to act on a prediction reasonably wants to know
which part of the molecule drove it. This is both a trust requirement and a
design one: knowing that a particular substructure is responsible for
predicted toxicity suggests what to change.

Attribution methods for graphs assign importance to nodes, edges or
subgraphs. Gradient-based approaches propagate sensitivity back to atom
features; perturbation-based approaches delete substructures and measure the
effect; and subgraph-search methods look for the minimal fragment that
preserves the prediction.

Two difficulties recur. First, molecular attributions are frequently unstable
--- small perturbations to the input, or retraining with a different seed,
produce materially different attributions, which is hard to reconcile with the
claim that they reveal something about the chemistry. Second, and more
fundamentally, deleting a substructure from a molecule does not produce
another molecule; it produces an invalid graph with unfilled valences. The
counterfactual the method requires is not chemically meaningful, and what the
model reports about it is not obviously interpretable.

The approaches that survive contact with practice tend to be those that stay
within chemical space --- comparing matched molecular pairs that differ by a
single well-defined substitution, for example, which is a manipulation
chemists already use and understand. This is a general pattern: attribution is
most reliable when the counterfactual is one the domain admits.

\section{Practical considerations}
\label{sec:practical}

Several details determine whether a graph model works in practice and receive
less attention than architecture.

\paragraph{Featurisation.}
The initial atom and bond features matter a great deal. A typical featuriser
supplies element, degree, formal charge, hybridisation, aromaticity, ring
membership and hydrogen count. As noted in Section~\ref{sec:expressive},
several of these encode information the message passing could not otherwise
derive, so featurisation choices interact with expressiveness claims. Papers
comparing architectures on different featurisers are not comparing
architectures.

\paragraph{Hydrogen treatment.}
Most implementations suppress explicit hydrogens and encode counts as atom
features, which shrinks graphs by roughly half. For properties dominated by
hydrogen bonding this loses relevant structure, and explicit-hydrogen models
sometimes do better at a substantial cost in graph size.

\paragraph{Tautomers and protonation.}
A molecule as drawn is one of several tautomeric forms, and the form present
at physiological pH may not be the one in the dataset. Different databases
standardise differently, so the same compound can appear as different graphs
across datasets. This is a persistent and under-reported source of label noise
in aggregated benchmarks.

\paragraph{Duplicate and near-duplicate leakage.}
Public datasets aggregated from multiple sources contain the same compound
more than once, sometimes with inconsistent labels. Where a random split
places duplicates on both sides, measured performance is inflated. Careful
deduplication by canonical structure is a precondition for meaningful
evaluation and is not always performed.

\section{Active learning and the discovery loop}
\label{sec:activelearning}

The deployment setting for these models is rarely a single prediction. It is a
loop: predict over a library, select compounds, test them, retrain on the
results, repeat. This changes what makes a model good.

In a one-shot setting the objective is accuracy over the whole distribution.
In a loop it is the quality of the selected batch, and those differ. A model
that is accurate on average but confidently wrong at the top of its ranking is
worse than a less accurate model that knows where it is uncertain, because
only the top of the ranking is ever tested. This is why the calibration
discussion of Section~\ref{sec:uncertainty} is not a refinement but a central
concern.

Acquisition strategy also matters more than architecture in early rounds.
Selecting purely by predicted activity exploits the current model and gathers
data where it is already confident, which teaches it little. Selecting by
uncertainty explores but wastes tests on compounds unlikely to be active.
Batch diversity constraints --- requiring selected compounds to be
structurally dissimilar from one another --- are often more valuable than
either, since a batch of near-identical molecules returns essentially one
measurement.

\paragraph{Evaluation implications.}
None of this is captured by a static benchmark. A model's value in a discovery
loop depends on how its errors are distributed relative to its confidence, and
on how the data it induces improves the next round. Simulating the loop
retrospectively over a screening library is possible and rarely done; it would
be a substantially more informative benchmark than the static splits currently
used, and would likely reorder the rankings again.

\section{Discussion}
\label{sec:discussion}

The theoretical literature on graph networks is unusually well developed.
We know the expressiveness ceiling exactly, we understand over-smoothing as
diffusion and over-squashing as capacity, and architectures exist that address
each. It is therefore striking how little of this translates into reliable
gains on the tasks the field cares about.

Our reading is that the binding constraint is elsewhere. The datasets are
small, biased toward explored chemistry, and evaluated in ways that reward
interpolation. In that regime the difference between a WL-limited and a
WL-exceeding architecture is not what determines performance; the difference
between a random and a scaffold split is far larger than the difference
between any two models under either.

This suggests a reallocation of effort. Better benchmarks --- prospective,
scaffold-split, with honest uncertainty quantification --- would do more for
the field than better aggregation functions. So would investment in data:
the successes in Section~\ref{sec:success} came from large consistent
screening sets, not from architectural novelty.

\section{Conclusion}
\label{sec:conclusion}

Message passing gives a principled way to learn functions of molecules,
respecting permutation symmetry by construction and adapting the
representation to the task. Its limits are well characterised: bounded by
1-WL, shallow because of over-smoothing, and unable to carry long-range signal
across bottlenecks. Geometric and transformer variants address specific
failures at specific costs.

What the field has not established is that these architectural distinctions
determine practical performance. Tuned fingerprint baselines remain
competitive, scaffold splits reorder published rankings, and the clearest
successes came from problems with abundant consistent data used as filters
rather than oracles. Until evaluation reflects the deployment setting, further
architectural refinement is likely to yield benchmark improvements that do not
transfer.

\begin{thebibliography}{99}

\bibitem{alon2021}
U.~Alon and E.~Yahav.
\newblock On the bottleneck of graph neural networks and its practical
  implications.
\newblock In \emph{ICLR}, 2021.

\bibitem{duvenaud2015}
D.~Duvenaud, D.~Maclaurin, J.~Aguilera-Iparraguirre, R.~Gómez-Bombarelli,
  T.~Hirzel, A.~Aspuru-Guzik, and R.~P. Adams.
\newblock Convolutional networks on graphs for learning molecular
  fingerprints.
\newblock In \emph{NeurIPS}, 2015.

\bibitem{gasteiger2020}
J.~Gasteiger, J.~Groß, and S.~Günnemann.
\newblock Directional message passing for molecular graphs.
\newblock In \emph{ICLR}, 2020.

\bibitem{gilmer2017}
J.~Gilmer, S.~S. Schoenholz, P.~F. Riley, O.~Vinyals, and G.~E. Dahl.
\newblock Neural message passing for quantum chemistry.
\newblock In \emph{ICML}, 2017.

\bibitem{hamilton2017}
W.~L. Hamilton, R.~Ying, and J.~Leskovec.
\newblock Inductive representation learning on large graphs.
\newblock In \emph{NeurIPS}, 2017.

\bibitem{jumper2021}
J.~Jumper et~al.
\newblock Highly accurate protein structure prediction with AlphaFold.
\newblock \emph{Nature}, 596:583--589, 2021.

\bibitem{kipf2017}
T.~N. Kipf and M.~Welling.
\newblock Semi-supervised classification with graph convolutional networks.
\newblock In \emph{ICLR}, 2017.

\bibitem{li2018}
Q.~Li, Z.~Han, and X.-M. Wu.
\newblock Deeper insights into graph convolutional networks for
  semi-supervised learning.
\newblock In \emph{AAAI}, 2018.

\bibitem{morris2019}
C.~Morris, M.~Ritzert, M.~Fey, W.~L. Hamilton, J.~E. Lenssen, G.~Rattan, and
  M.~Grohe.
\newblock Weisfeiler and Leman go neural: Higher-order graph neural networks.
\newblock In \emph{AAAI}, 2019.

\bibitem{rampasek2022}
L.~Rampášek, M.~Galkin, V.~P. Dwivedi, A.~T. Luu, G.~Wolf, and D.~Beaini.
\newblock Recipe for a general, powerful, scalable graph transformer.
\newblock In \emph{NeurIPS}, 2022.

\bibitem{satorras2021}
V.~G. Satorras, E.~Hoogeboom, and M.~Welling.
\newblock E($n$) equivariant graph neural networks.
\newblock In \emph{ICML}, 2021.

\bibitem{schutt2017}
K.~T. Schütt, P.-J. Kindermans, H.~E. Sauceda, S.~Chmiela, A.~Tkatchenko, and
  K.-R. Müller.
\newblock SchNet: A continuous-filter convolutional neural network for
  modeling quantum interactions.
\newblock In \emph{NeurIPS}, 2017.

\bibitem{stokes2020}
J.~M. Stokes et~al.
\newblock A deep learning approach to antibiotic discovery.
\newblock \emph{Cell}, 180(4):688--702, 2020.

\bibitem{velickovic2018}
P.~Veličković, G.~Cucurull, A.~Casanova, A.~Romero, P.~Liò, and Y.~Bengio.
\newblock Graph attention networks.
\newblock In \emph{ICLR}, 2018.

\bibitem{wu2018}
Z.~Wu, B.~Ramsundar, E.~N. Feinberg, J.~Gomes, C.~Geniesse, A.~S. Pappu,
  K.~Leswing, and V.~Pande.
\newblock MoleculeNet: A benchmark for molecular machine learning.
\newblock \emph{Chemical Science}, 9:513--530, 2018.

\bibitem{xu2019}
K.~Xu, W.~Hu, J.~Leskovec, and S.~Jegelka.
\newblock How powerful are graph neural networks?
\newblock In \emph{ICLR}, 2019.

\bibitem{ying2021}
C.~Ying, T.~Cai, S.~Luo, S.~Zheng, G.~Ke, D.~He, Y.~Shen, and T.-Y. Liu.
\newblock Do transformers really perform badly for graph representation?
\newblock In \emph{NeurIPS}, 2021.

\end{thebibliography}

\end{document}
